\uuid{jq7e}
\exo7id{800}
\titre{exo7 800}
\auteur{cousquer}
\organisation{exo7}
\datecreate{2003-10-01}
\isIndication{false}
\isCorrection{false}
\chapitre{Calcul d'intégrales}
\sousChapitre{Théorie}
\cpge{mpsi-mp2i.fonctions-valeurs-reelles-complexes.a, pcsi.techniques-calcul-fonctions.a, ptsi.techniques-calcul-fonctions.a, mpsi-mp2i.limites-continuite.b, pcsi.limites-continuite.b, ptsi.limites-continuite.b, mpsi-mp2i.integration.c, pcsi.integration.b, ptsi.integration.b, bcpst1.limites-continuite.c, bcpst1.integration-segment.a, bcpst1.integration-segment.d}
\module{Analyse}
\niveau{L1}
\difficulte{}

\contenu{
\question{
Soit $f$ une fonction continue et strictement croissante de $[a,b]$ sur
$[\alpha,\beta]$. On note $g$ la fonction réciproque de~$f$. Montrer que
$$\int_a^bf(x)\,dx+\int_\alpha^\beta g(x)\,dx = b\beta-a\alpha$$
}
}
